\scanpage{9}
\mspage{8}

\deleted{\uline{Proposition}}

\deleted{\uline{Définition} $\overbrace{\text{Segment, intervalle, demi-intervalle}}^{\inserted{\text{tronçons}}}$
(comme structure de réseau).}
\deleted{Segment}

\uline{Définition}\quad Un ens. ordonné \inserted{$\mathcal L$} est appelé un \uline{tronçon} \uline{ordonné} si

\unsure{a}) $\forall\,a\in\mathcal L$ tel que $\mathcal L_{>a}\deleted{\unsure{\sim|=}}\unsure{\ne|=}\varnothing$, $\mathcal L_{>a}$ est filtrant décroissant
\deleted{i.e. $\forall\,x,y$ $\unsure{a|x},y$ tels que} \inserted{si $a,y,\unsure{a|x}\in\mathcal L$ tels que}
$a>a$, $y>a$ alors $\exists\,z\in\mathcal L$ avec $x,y\geq z>a$

$\deleted{\unsure{b|a}}a'$) Axiome dual : si $a,x,y\in\mathcal L$ tels q
$\unsure{x|a}<a$, $y<a$ alors $\exists\,z\in\mathcal L$ avec $a,y\leq z<a$

b) $\mathcal L$ est filt\unsure{ant|rant} croissant et filtrant décroissant i.e.
$x,y\in\mathcal L\Rightarrow\exists\,a,b\in\mathcal L$ avec $a\leq x,y\leq b$

c) Si $x,y\in\mathcal L$ tels que $x<y$, $\exists\,z$ tel que $x<z<y$ .

d) $\deleted{\mathcal L\,\unsure{\ne|=}}$ $\operatorname{Card}\mathcal L>1$ \deleted{\unsure{(imp.|(}}

% slanted
\margin{\uline{NB} les cond $a\ a'\ b)\ c)$ sont satisfaites si $\operatorname{card}\mathcal L\leq1$ !}

On dit que $\mathcal L$ est un \uline{segment} si de plus $\mathcal L$ admet un plus petit et
\deleted{\unsure{(p.|(}} un plus grand élément, un \uline{intervalle} s'il n'admet ni l'un ni l'autre,

% paper watermark, visible from the reverse
\rotatebox{90}{\textsc{\unsure{n}g Extra Stro}}
